\section{Introduction Générale}

\subsection{Contexte général}

Dans un secteur des services numériques de plus en plus concurrentiel, la capacité d'une entreprise à
positionner correctement son offre commerciale constitue un avantage stratégique déterminant. Les
entreprises de développement web et mobile évoluent dans un marché où les prix pratiqués, la qualité
perçue des prestations et la réputation des concurrents évoluent en permanence. Or, la collecte et
l'exploitation de ces informations reposent encore, dans la majorité des cas, sur des démarches manuelles,
ponctuelles et peu structurées.

Le développement récent des techniques d'intelligence artificielle générative, et en particulier des
architectures de \textit{Retrieval-Augmented Generation} (RAG), ouvre de nouvelles perspectives pour
automatiser la collecte, la structuration et l'exploitation de données non structurées issues du web. Combinées
à des techniques de traitement automatique du langage naturel (NLP) telles que l'analyse de sentiment,
ces approches permettent de transformer des données brutes --- prix affichés, avis clients --- en
informations directement exploitables par les équipes commerciales.

C'est dans ce contexte que s'inscrit le stage de fin d'études réalisé au sein d'ICHTUS IT, entreprise
spécialisée dans le développement de plateformes web et d'applications mobiles.

\subsection{Problématique et enjeux}

Aujourd'hui, ICHTUS IT ne dispose d'aucun outil centralisant les données tarifaires et les avis clients de
ses concurrents. Chaque commercial doit effectuer ses propres recherches ponctuelles, sans historique ni
vision d'ensemble sur le positionnement du marché. Cette absence d'outillage ralentit le chiffrage des
offres et prive l'équipe commerciale d'arguments de vente construits sur des données réelles.

La problématique de ce mémoire peut donc être formulée ainsi : \textit{comment automatiser la collecte des
données tarifaires et des avis clients des entreprises concurrentes, et les rendre exploitables en temps réel
via un agent conversationnel, pour permettre à l'équipe commerciale de prendre des décisions de pricing et
de prospection plus rapides et mieux argumentées ?}

Cette problématique soulève des enjeux à plusieurs niveaux :
\begin{itemize}
    \item \textbf{Enjeux scientifiques et techniques} : concevoir un pipeline robuste de collecte de données
    web non structurées, une architecture de recherche vectorielle performante, et un système de
    classification de sentiment fiable en l'absence de données annotées.
    \item \textbf{Enjeux industriels} : réduire le temps de recherche manuelle des commerciaux et fiabiliser
    la prise de décision en matière de pricing.
    \item \textbf{Enjeux socio-économiques} : renforcer la compétitivité d'ICHTUS IT sur son marché grâce à
    une meilleure connaissance de son environnement concurrentiel.
\end{itemize}

\subsection{Objectifs du stage et du mémoire}

Ce stage vise à concevoir et développer un système capable de :
\begin{itemize}
    \item collecter automatiquement les prix pratiqués par les entreprises concurrentes du secteur du
    développement web et mobile ;
    \item mettre ces données à disposition de l'équipe commerciale via un agent conversationnel interrogeable
    en langage naturel ;
    \item analyser les avis clients laissés sur les concurrents afin d'en extraire automatiquement le
    sentiment global ainsi que les points forts et points faibles récurrents.
\end{itemize}

L'objectif du mémoire est, quant à lui, de documenter la démarche scientifique et technique ayant permis
d'aboutir à ce système, de la revue de littérature jusqu'à l'évaluation des résultats obtenus.

\subsection{Hypothèses et questions de recherche}

Ce travail repose sur les hypothèses suivantes :
\begin{itemize}
    \item il est possible de collecter de manière automatisée et fiable des données tarifaires et des avis
    clients publiés en ligne par des entreprises concurrentes, malgré l'hétérogénéité de leur présentation ;
    \item une architecture RAG appuyée sur une recherche par similarité vectorielle permet de restituer une
    information de marché pertinente et à jour, sans nécessiter de réentraînement du modèle de langage ;
    \item une analyse de sentiment en \textit{zero-shot}, sans fine-tuning ni jeu de données annoté, offre une
    précision suffisante pour être exploitée par des utilisateurs non techniques.
\end{itemize}

Ces hypothèses guident les questions de recherche suivantes : dans quelle mesure un pipeline ETL appliqué à
des données web non structurées permet-il d'alimenter un système d'aide à la décision fiable ? Et comment
combiner recherche vectorielle et extraction de sentiment pour produire une réponse à la fois quantitative
et qualitative sur le positionnement des concurrents ?

\subsection{Méthodologie générale adoptée}

La démarche suivie s'articule en trois grandes phases correspondant aux trois couches du système développé :

\begin{enumerate}
    \item une phase de \textbf{collecte de données}, reposant sur un scraping intelligent automatisé
    (ScrapeGraphAI) des pages de tarifs et d'avis clients, suivie d'un pipeline de transformation
    (parsing, structuration) inspiré des démarches ETL classiques ;
    \item une phase de \textbf{stockage et de traitement}, avec une base de données Supabase/PostgreSQL
    indexée pour la recherche sémantique grâce à l'extension \texttt{pgvector} ;
    \item une phase d'\textbf{intelligence artificielle}, combinant un système de RAG pour l'interrogation en
    langage naturel du marché, et un module d'analyse de sentiment basé sur un grand modèle de langage en
    extraction \textit{zero-shot}.
\end{enumerate}

Cette méthodologie est détaillée et justifiée dans les chapitres consacrés à l'état de l'art et à la
méthodologie du présent mémoire.

\subsection{Organisation du mémoire}

Ce mémoire s'organise de la manière suivante. Le chapitre 2 présente l'établissement académique
d'origine, l'INSI. Le chapitre 3 présente l'entreprise d'accueil, ICHTUS IT. Le chapitre 4 expose le
cadrage du projet, les besoins identifiés et les contraintes rencontrées. Le chapitre 5 propose une revue
de la littérature scientifique et technique relative au sujet traité. Le chapitre 6 détaille la méthodologie
adoptée. Les chapitres 7 et 8 présentent respectivement l'exploration et la préparation des données, puis
la modélisation et les expérimentations menées. Le chapitre 9 décrit le déploiement du système sous forme
de démonstrateur. Le chapitre 10 propose une discussion critique des résultats obtenus, avant que le
chapitre 11 n'expose les limites du travail et les perspectives envisagées. Le mémoire se conclut par une
synthèse générale des apports scientifiques, techniques et professionnels de ce stage.